%==============================================================================
% APPENDIX A: Julia Quickstart Reference
%==============================================================================
\chapter{Julia Quickstart Reference}
\label{app:julia}

\begin{abstract}
This appendix provides a rapid-reference guide to Julia syntax for scientists and engineers. Every chapter in this book includes Julia code; this appendix ensures you can follow along even if you are new to the language.
\end{abstract}

%------------------------------------------------------------------------------
\section{Basic Syntax}
\label{app:julia:syntax}

\begin{description}
    \item[Variables] Assign with \texttt{=}. Types are inferred:
    \begin{lstlisting}[language=Julia]
x = 3.14          # Float64
n = 42            # Int64
s = "hello"       # String
b = true          # Bool
    \end{lstlisting}
    \item[Comments] Use \texttt{\#} for single-line comments.
    \item[Semicolons] \texttt{;} suppresses output in REPL.
\end{description}

%------------------------------------------------------------------------------
\section{Arrays and Matrices}
\label{app:julia:arrays}

\begin{lstlisting}[language=Julia]
# Vectors
v = [1, 2, 3]              # Column vector
w = [1.0, 2.0, 3.0]        # Float vector

# Matrices
A = [1 2 3; 4 5 6; 7 8 9]  # 3x3 matrix
I = Matrix(Identity, 3)     # 3x3 identity

# Operations
B = A'                     # Transpose
C = A * v                  # Matrix-vector product
D = A \ v                  # Solve Ax = v

# Element-wise operations
x = w .^ 2                 # Square each element
y = log.(v)                # Element-wise log
    \end{lstlisting}

\section{Functions}
\label{app:julia:functions}

\begin{lstlisting}[language=Julia]
# Function definition
function f(x::Float64)
    return x^2 + 1
end

# One-liner
g(x) = x^2 + 1

# Anonymous function
h = x -> x^2 + 1

# Multiple dispatch
foo(x::Int) = "integer"
foo(x::String) = "string"
    \end{lstlisting}

\section{Loops and Comprehensions}
\label{app:julia:loops}

\begin{lstlisting}[language=Julia]
# For loop
for i in 1:10
    println(i)
end

# While loop
i = 1
while i < 10
    i += 1
end

# Array comprehension
squares = [i^2 for i in 1:10]

# Filtered comprehension
evens = [i for i in 1:20 if iseven(i)]
    \end{lstlisting}

\section{Linear Algebra}
\label{app:julia:linalg}

\begin{lstlisting}[language=Julia]
using LinearAlgebra

A = rand(3, 3)
b = rand(3)

# Factorization
F = lu(A)
x = F \ b

# Eigenvalues
vals, vecs = eigvals(A), eigvecs(A)

# SVD
U, S, V = svd(A)

# Sparse matrices
using SparseArrays
S = sprand(100, 100, 0.01)
    \end{lstlisting}

\section{Plotting}
\label{app:julia:plotting}

\begin{lstlisting}[language=Julia]
using Plots

# Basic plot
x = linspace(0, 2pi, 100)
plot(x, sin.(x), label="sin(x)")
scatter!(x, cos.(x), label="cos(x)")
savefig("plot.png")

# 3D surface
using PyPlot
x = y = linspace(-2, 2, 50)
[X, Y] = meshgrid(x, y)
Z = X.*exp.(-X.^2 - Y.^2)
surface(X, Y, Z)
    \end{lstlisting}

\section{Differential Equations}
\label{app:julia:ode}

\begin{lstlisting}[language=Julia]
using DifferentialEquations

# Lorenz system
function lorenz(u, p, t)
    sigma, rho, beta = p
    du1 = sigma * (u2 - u1)
    du2 = u1 * (rho - u3) - u2
    du3 = u1 * u2 - beta * u3
    return [du1, du2, du3]
end

u0 = [1.0, 0.0, 0.0]
p = [10.0, 28.0, 8/3]
tspan = (0.0, 100.0)
prob = ODEProblem(lorenz, u0, tspan, p)
sol = solve(prob, Tsit5())
    \end{lstlisting}

\section{Package Management}
\label{app:julia:pkg}

\begin{lstlisting}[language=Julia]
using Pkg

# Add packages
Pkg.add("LinearAlgebra")
Pkg.add("Plots")
Pkg.add("DifferentialEquations")

# Instantiate environment
Pkg.instantiate()

# Update all packages
Pkg.update()
    \end{lstlisting}

%------------------------------------------------------------------------------
\section{Summary}
\label{app:julia:summary}

Key Julia features for scientific computing:
\begin{enumerate}
    \item \textbf{Multiple dispatch}: Functions specialize on argument types.
    \item \textbf{Type stability}: Write code that the compiler can optimize.
    \item \textbf{Arrays are 1-indexed}: Unlike Python/MATLAB.
    \item \textbf{Vectorized operations}: Use dot notation (\texttt{.}) for element-wise ops.
    \item \textbf{Unicode support}: \texttt{x\_1}, \texttt{\\alpha} work in identifiers.
\end{enumerate}

\index{Julia, quickstart}
\index{packages, management}
\index{differential equations, Julia}
\index{linear algebra, Julia}
\index{plotting, Julia}
